\documentclass[../main.tex]{subfiles}

\begin{document}
\subsection{Numerical Methods for Ordinary Differential Equations}
\label{subsec:ode-methods}


\textBox{%
An ordinary differential equation specifies the instantaneous velocity of a continuous state trajectory. A numerical time-stepping method replaces this continuous flow by a sequence of computable discrete states. The core mathematical questions in numerical ODE integration are consistency, zero-stability, global convergence, stiffness, and geometric structure preservation.
}

\subsubsection{Initial Value Problems and Well-Posedness}

\addDef{Initial Value Problem}{%
An initial value problem (IVP) for a system of ordinary differential equations has the form
\begin{align*}
    \by'(t) = \boldsymbol{f}(t, \by(t)), \qquad \by(t_0) = \by_0,
\end{align*}
where $t \in [t_0, T]$ and $\by(t) \in \fR^m$. The function $\boldsymbol{f}: [t_0, T] \times \fR^m \to \fR^m$ is the vector field.
}
\label{def:ode-ivp}

\addThm{Picard--Lindelöf Existence and Uniqueness}{%
Let $\boldsymbol{f}(t, \by)$ be continuous in $t$ and uniformly Lipschitz continuous in $\by$ on the domain $[t_0, T] \times \fR^m$:
\begin{align*}
    \|\boldsymbol{f}(t, \bu) - \boldsymbol{f}(t, \bv)\| \leq L \|\bu - \bv\| \qquad \text{for all } \bu, \bv \in \fR^m,
\end{align*}
where $L < \infty$ is the Lipschitz constant. Then the IVP has a unique continuous solution $\by(t)$ on $[t_0, T]$.
}
\label{thm:ode-existence-uniqueness}

\addPrf{%
The differential equation $\by' = \boldsymbol{f}(t, \by)$ with initial condition $\by(t_0) = \by_0$ is equivalent to the Picard integral equation:
\begin{align*}
    \by(t) = \by_0 + \int_{t_0}^t \boldsymbol{f}(s, \by(s))\,ds.
\end{align*}
Assume first that $L>0$. Define the Picard operator $\mathcal{T}$ on the Banach space $C([t_0, T]; \fR^m)$ equipped with the weighted norm $\|\by\|_* = \max_{t \in [t_0, T]} e^{-2L(t-t_0)} \|\by(t)\|$.
For any two functions $\bu, \bv$:
\begin{align*}
    \|(\mathcal{T}\bu)(t) - (\mathcal{T}\bv)(t)\| &\leq \int_{t_0}^t \|\boldsymbol{f}(s, \bu(s)) - \boldsymbol{f}(s, \bv(s))\|\,ds \\
    &\leq L \int_{t_0}^t \|\bu(s) - \bv(s)\|\,ds \\
    &= L \int_{t_0}^t e^{2L(s-t_0)} \left( e^{-2L(s-t_0)} \|\bu(s) - \bv(s)\| \right) ds \\
    &\leq L \|\bu - \bv\|_* \int_{t_0}^t e^{2L(s-t_0)}\,ds \\
    &= \frac{1}{2} \left(e^{2L(t-t_0)}-1\right)\|\bu-\bv\|_*.
\end{align*}
Multiplying through by $e^{-2L(t-t_0)}$ yields $\|\mathcal{T}\bu - \mathcal{T}\bv\|_* \leq \frac{1}{2} \|\bu - \bv\|_*$.
Because the operator $\mathcal{T}$ is a strict contraction with modulus $1/2$ on a complete metric space, the Banach Fixed-Point Theorem guarantees the existence of a unique fixed point $\by^* = \mathcal{T}\by^*$, which is the unique solution to the IVP.
If $L=0$, then $\boldsymbol f(t,\bu)=\boldsymbol f(t,\bv)$ for every $\bu,\bv$, so $\mathcal{T}$ has contraction modulus zero in the ordinary supremum norm.
}

\subsubsection{Consistency, Stability, and the Discrete Gronwall Lemma}

\textBox{%
A numerical time-stepping method is convergent if the computed trajectory approaches the exact trajectory as the step size $h \to 0$. We quantify error at each step and track its accumulation over a finite time interval.
}

\addDef{Local Truncation Error and One-Step Defect}{%
Let $\by_{n+1} = \Phi_h(t_n, \by_n)$ denote an autonomous or non-autonomous one-step method. The local truncation error is the defect of the exact solution:
\begin{align*}
    \boldsymbol{\tau}_{n+1} = \frac{\by(t_{n+1}) - \Phi_h(t_n, \by(t_n))}{h}.
\end{align*}
The method is \textbf{consistent of order $p$} if $\|\boldsymbol{\tau}_{n+1}\| = O(h^p)$ as $h \to 0$ for all sufficiently smooth solutions.
}
\label{def:ode-local-error}

\addLem{Discrete Gronwall Lemma}{%
Let $\{e_n\}_{n=0}^N$ be a sequence of non-negative real numbers satisfying the linear recurrence
\begin{align*}
    e_{n+1} \leq (1 + \alpha) e_n + \beta, \qquad n = 0, 1, \dots, N-1,
\end{align*}
where $\alpha, \beta \geq 0$. For $\alpha>0$,
\begin{align*}
    e_n \leq e^{n \alpha} e_0 + \beta \frac{e^{n \alpha} - 1}{\alpha}.
\end{align*}
For $\alpha=0$, the corresponding bound is $e_n\leq e_0+n\beta$.
}
\label{lem:discrete-gronwall}

\addPrf{%
Expanding the recurrence by induction gives
\begin{align*}
    e_n \leq (1 + \alpha)^n e_0 + \beta \sum_{k=0}^{n-1} (1 + \alpha)^k = (1 + \alpha)^n e_0 + \beta \frac{(1 + \alpha)^n - 1}{\alpha}.
\end{align*}
Using the standard inequality $1 + \alpha \leq e^\alpha$ for all $\alpha \geq 0$, we have $(1 + \alpha)^n \leq (e^\alpha)^n = e^{n \alpha}$, which completes the proof.
}

\addThm{Fundamental Theorem of One-Step Convergence}{%
Let $\Phi_h$ be a consistent one-step method of order $p$ applied to an IVP whose vector field $\boldsymbol{f}$ is Lipschitz continuous with constant $L$. Suppose the numerical increment function satisfies $\|\Phi_h(t, \bu) - \Phi_h(t, \bv)\| \leq (1 + \Lambda h) \|\bu - \bv\|$ for some constant $\Lambda \geq 0$.
Then on the fixed interval $[t_0, T]$ with $Nh = T - t_0$, the global error $\be_n = \by(t_n) - \by_n$ satisfies
\begin{align*}
    \max_{0 \leq n \leq N} \|\be_n\| \leq e^{\Lambda(T - t_0)} \|\be_0\| + \frac{e^{\Lambda(T - t_0)} - 1}{\Lambda} \max_k \|\boldsymbol{\tau}_k\|,
\end{align*}
where the quotient is interpreted as $T-t_0$ when $\Lambda=0$.
Consequently, when the numerical and exact solutions use the same initial value, a local error of order $O(h^{p+1})$ yields a global error of order $O(h^p)$.
}
\label{thm:local-to-global-ode}

\addPrf{%
Express the error at step $n+1$ by subtracting the numerical update from the exact solution:
\begin{align*}
    \be_{n+1} &= \by(t_{n+1}) - \by_{n+1} \\
    &= \left[ \Phi_h(t_n, \by(t_n)) + h \boldsymbol{\tau}_{n+1} \right] - \Phi_h(t_n, \by_n) \\
    &= \left[ \Phi_h(t_n, \by(t_n)) - \Phi_h(t_n, \by_n) \right] + h \boldsymbol{\tau}_{n+1}.
\end{align*}
Taking norms and applying the triangle inequality yields
\begin{align*}
    \|\be_{n+1}\| \leq (1 + \Lambda h) \|\be_n\| + h \max_k \|\boldsymbol{\tau}_k\|.
\end{align*}
Applying the Discrete Gronwall Lemma (Lem~\ref{lem:discrete-gronwall}) with $\alpha = \Lambda h$ and $\beta = h \max_k \|\boldsymbol{\tau}_k\|$ gives
\begin{align*}
    \|\be_n\| \leq e^{n \Lambda h} \|\be_0\| + h \max_k \|\boldsymbol{\tau}_k\| \frac{e^{n \Lambda h} - 1}{\Lambda h} = e^{\Lambda(t_n - t_0)} \|\be_0\| + \frac{e^{\Lambda(t_n - t_0)} - 1}{\Lambda} \max_k \|\boldsymbol{\tau}_k\|.
\end{align*}
Since $t_n - t_0 \leq T - t_0$ and $\max_k \|\boldsymbol{\tau}_k\| = O(h^p)$, the global error is bounded by $O(h^p)$ when initial data are exact ($\be_0 = \bzero$).
}

\subsubsection{Euler Methods and Stability Functions}

\addDef{Forward Euler, Backward Euler, and Trapezoidal Methods}{%
For time step $h = t_{n+1} - t_n$:
\begin{enumerate}
    \item \textbf{Forward Euler (explicit):} $\by_{n+1} = \by_n + h \boldsymbol{f}(t_n, \by_n)$.
    \item \textbf{Backward Euler (implicit):} $\by_{n+1} = \by_n + h \boldsymbol{f}(t_{n+1}, \by_{n+1})$.
    \item \textbf{Trapezoidal Rule (implicit Crank--Nicolson):}
    \begin{align*}
        \by_{n+1} = \by_n + \frac{h}{2} \left[ \boldsymbol{f}(t_n, \by_n) + \boldsymbol{f}(t_{n+1}, \by_{n+1}) \right].
    \end{align*}
\end{enumerate}
}
\label{def:euler-methods}

\addThm{Stability Functions for Linear Test Equations}{%
Applied to the Dahlquist test equation $y' = \lambda y$ with $z = h\lambda \in \fC$, the numerical update satisfies $y_{n+1} = R(z) y_n$. The stability functions are:
\begin{align*}
    R_{\mathrm{FE}}(z) = 1 + z, \qquad
    R_{\mathrm{BE}}(z) = \frac{1}{1 - z}, \qquad
    R_{\mathrm{TR}}(z) = \frac{1 + z/2}{1 - z/2}.
\end{align*}
The region of absolute stability is $\mathcal{S} = \{z \in \fC : |R(z)| \leq 1\}$.
Backward Euler and the Trapezoidal Rule are A-stable: $\fC^- = \{z \in \fC : \operatorname{Re}(z) \leq 0\} \subseteq \mathcal{S}$.
Backward Euler is additionally L-stable: $\lim_{z \to -\infty} R_{\mathrm{BE}}(z) = 0$.
}
\label{thm:euler-stability-functions}

\addPrf{%
For Forward Euler: $y_{n+1} = y_n + h\lambda y_n = (1 + z)y_n$.
For Backward Euler: $y_{n+1} = y_n + z y_{n+1} \implies (1 - z)y_{n+1} = y_n \implies y_{n+1} = \frac{1}{1 - z} y_n$.
For Trapezoidal Rule: $y_{n+1} = y_n + \frac{z}{2}(y_n + y_{n+1}) \implies (1 - z/2)y_{n+1} = (1 + z/2)y_n \implies y_{n+1} = \frac{1 + z/2}{1 - z/2} y_n$.
Let $z = x + iy$ with $x \leq 0$.
For Backward Euler, $|1 - z|^2 = (1 - x)^2 + y^2 \geq 1 + y^2 \geq 1$, so $|R_{\mathrm{BE}}(z)| = \frac{1}{|1 - z|} \leq 1$.
For Trapezoidal Rule,
\begin{align*}
    |1 + z/2|^2 = (1 + x/2)^2 + (y/2)^2 = 1 + x + \frac{x^2 + y^2}{4}, \\
    |1 - z/2|^2 = (1 - x/2)^2 + (y/2)^2 = 1 - x + \frac{x^2 + y^2}{4}.
\end{align*}
Because $x \leq 0$, we have $1 + x \leq 1 - x$, which implies $|1 + z/2|^2 \leq |1 - z/2|^2$ and $|R_{\mathrm{TR}}(z)| \leq 1$.
Thus both methods are A-stable.
Along the negative real axis, $\lim_{x \to -\infty} |R_{\mathrm{BE}}(x)| = \lim_{x \to -\infty} \frac{1}{1 - x} = 0$, proving L-stability. In contrast, $\lim_{x \to -\infty} R_{\mathrm{TR}}(x) = -1$, so the Trapezoidal Rule is not L-stable.
}

\subsubsection{Runge--Kutta Methods}

\textBox{%
Runge--Kutta methods achieve high orders of accuracy without storing historical steps by evaluating intermediate stages within each time step.
}

\addDef{Classical Fourth-Order Runge--Kutta Method (RK4)}{%
For an autonomous ODE $\by' = \boldsymbol{f}(\by)$, the classical RK4 method computes four stages:
\begin{align*}
    \boldsymbol{k}_1 &= \boldsymbol{f}(\by_n), \\
    \boldsymbol{k}_2 &= \boldsymbol{f}\left( \by_n + \frac{h}{2}\boldsymbol{k}_1 \right), \\
    \boldsymbol{k}_3 &= \boldsymbol{f}\left( \by_n + \frac{h}{2}\boldsymbol{k}_2 \right), \\
    \boldsymbol{k}_4 &= \boldsymbol{f}\left( \by_n + h\boldsymbol{k}_3 \right),
\end{align*}
and sets the update
\begin{align*}
    \by_{n+1} = \by_n + \frac{h}{6} \left( \boldsymbol{k}_1 + 2\boldsymbol{k}_2 + 2\boldsymbol{k}_3 + \boldsymbol{k}_4 \right).
\end{align*}
The one-step defect is $O(h^5)$. With the normalization in Def~\ref{def:ode-local-error}, the local truncation error and the global error are both $O(h^4)$.
}
\label{def:rk4}

\addThm{RK4 Stability Polynomial}{%
Applied to the scalar linear test equation $y' = \lambda y$, the stability function of RK4 is
\begin{align*}
    R(z) = 1 + z + \frac{z^2}{2} + \frac{z^3}{6} + \frac{z^4}{24} = \sum_{k=0}^4 \frac{z^k}{k!}, \qquad z = h\lambda.
\end{align*}
}
\label{thm:rk4-stability}

\addPrf{%
Substituting $f(y) = \lambda y$ into the RK4 stages:
\begin{align*}
    k_1 &= \lambda y_n, \\
    k_2 &= \lambda \left( y_n + \frac{h}{2} k_1 \right) = \lambda \left( 1 + \frac{z}{2} \right) y_n, \\
    k_3 &= \lambda \left( y_n + \frac{h}{2} k_2 \right) = \lambda \left( 1 + \frac{z}{2} + \frac{z^2}{4} \right) y_n, \\
    k_4 &= \lambda \left( y_n + h k_3 \right) = \lambda \left( 1 + z + \frac{z^2}{2} + \frac{z^3}{4} \right) y_n.
\end{align*}
Substituting these into the update formula:
\begin{align*}
    y_{n+1} &= y_n + \frac{h}{6} \left[ \lambda + 2\lambda\left(1 + \frac{z}{2}\right) + 2\lambda\left(1 + \frac{z}{2} + \frac{z^2}{4}\right) + \lambda\left(1 + z + \frac{z^2}{2} + \frac{z^3}{4}\right) \right] y_n \\
    &= y_n + \frac{z}{6} \left[ 1 + (2 + z) + \left(2 + z + \frac{z^2}{2}\right) + \left(1 + z + \frac{z^2}{2} + \frac{z^3}{4}\right) \right] y_n \\
    &= y_n + \frac{z}{6} \left[ 6 + 3z + z^2 + \frac{z^3}{4} \right] y_n \\
    &= \left( 1 + z + \frac{z^2}{2} + \frac{z^3}{6} + \frac{z^4}{24} \right) y_n.
\end{align*}
}

\subsubsection{Symplectic Integrators and Geometric Mechanics}

\textBox{%
Hamiltonian systems preserve geometric phase-space volume and energy. Standard Runge--Kutta methods dissipate or amplify energy over long integration horizons. Symplectic integrators preserve the underlying canonical symplectic 2-form.
}

\addDef{Velocity Verlet Algorithm}{%
For a separable Hamiltonian $H(\bq, \bv) = \frac{1}{2}\|\bv\|_2^2 + V(\bq)$ with equations of motion $\bq' = \bv$ and $\bv' = \bF(\bq) = -\nabla V(\bq)$, the Velocity Verlet method is:
\begin{align*}
    \bv_{n+1/2} &= \bv_n + \frac{h}{2} \bF(\bq_n), \\
    \bq_{n+1} &= \bq_n + h \bv_{n+1/2}, \\
    \bv_{n+1} &= \bv_{n+1/2} + \frac{h}{2} \bF(\bq_{n+1}).
\end{align*}
}
\label{def:velocity-verlet}

\addThm{Symplecticity and Area Preservation of Velocity Verlet}{%
The one-step Velocity Verlet map $\Psi_h: (\bq_n, \bv_n) \mapsto (\bq_{n+1}, \bv_{n+1})$ preserves the canonical symplectic structure. In particular, the Jacobian matrix $\bJ = \frac{\partial(\bq_{n+1}, \bv_{n+1})}{\partial(\bq_n, \bv_n)}$ satisfies
\begin{align*}
    \det(\bJ) = 1.
\end{align*}
Equivalently, its Jacobian satisfies $\bJ^T\boldsymbol{\Omega}\bJ=\boldsymbol{\Omega}$ for the canonical symplectic matrix $\boldsymbol{\Omega}$. Symplecticity implies $\det(\bJ)=1$, so Velocity Verlet also preserves phase-space volume.
}
\label{thm:verlet-symplecticity}

\addPrf{%
Decompose the Velocity Verlet step into the composition of three elementary shears:
\begin{align*}
    \Psi_h = \psi_3 \circ \psi_2 \circ \psi_1,
\end{align*}
where:
\begin{align*}
    \psi_1(\bq, \bv) &= \left( \bq,\, \bv + \frac{h}{2}\bF(\bq) \right), \\
    \psi_2(\bq, \bv) &= (\bq + h\bv,\, \bv), \\
    \psi_3(\bq, \bv) &= \left( \bq,\, \bv + \frac{h}{2}\bF(\bq) \right).
\end{align*}
Evaluating the Jacobian of each shear map:
\begin{align*}
    \bJ_1 = \begin{pmatrix} \bI & \bzero \\ \frac{h}{2}\nabla\bF(\bq) & \bI \end{pmatrix}, \qquad
    \bJ_2 = \begin{pmatrix} \bI & h\bI \\ \bzero & \bI \end{pmatrix}, \qquad
    \bJ_3 = \begin{pmatrix} \bI & \bzero \\ \frac{h}{2}\nabla\bF(\bq_{n+1}) & \bI \end{pmatrix}.
\end{align*}
Because $\bF=-\nabla V$, the Hessian block $\bS = \frac{h}{2}\nabla\bF = -\frac{h}{2}\nabla^2V$ is symmetric ($\bS^T = \bS$).
With canonical symplectic matrix $\boldsymbol{\Omega} = \begin{pmatrix} \bzero & \bI \\ -\bI & \bzero \end{pmatrix}$, evaluating the quadratic form for $\bJ_1$:
\begin{align*}
    \bJ_1^T \boldsymbol{\Omega} \bJ_1 &= \begin{pmatrix} \bI & \bS \\ \bzero & \bI \end{pmatrix} \begin{pmatrix} \bzero & \bI \\ -\bI & \bzero \end{pmatrix} \begin{pmatrix} \bI & \bzero \\ \bS & \bI \end{pmatrix} \\
    &= \begin{pmatrix} -\bS & \bI \\ -\bI & \bzero \end{pmatrix} \begin{pmatrix} \bI & \bzero \\ \bS & \bI \end{pmatrix} = \begin{pmatrix} -\bS + \bS & \bI \\ -\bI & \bzero \end{pmatrix} = \boldsymbol{\Omega}.
\end{align*}
Similarly, for the translation shear $\bJ_2$:
\begin{align*}
    \bJ_2^T \boldsymbol{\Omega} \bJ_2 &= \begin{pmatrix} \bI & \bzero \\ h\bI & \bI \end{pmatrix} \begin{pmatrix} \bzero & \bI \\ -\bI & \bzero \end{pmatrix} \begin{pmatrix} \bI & h\bI \\ \bzero & \bI \end{pmatrix} \\
    &= \begin{pmatrix} \bzero & \bI \\ -\bI & h\bI \end{pmatrix} \begin{pmatrix} \bI & h\bI \\ \bzero & \bI \end{pmatrix} = \begin{pmatrix} \bzero & \bI \\ -\bI & -h\bI + h\bI \end{pmatrix} = \boldsymbol{\Omega}.
\end{align*}
The third shear $\bJ_3$ has the same form as $\bJ_1$ and satisfies $\bJ_3^T \boldsymbol{\Omega} \bJ_3 = \boldsymbol{\Omega}$.
By the chain rule, $\bJ = \bJ_3 \bJ_2 \bJ_1$, so
\begin{align*}
    \bJ^T \boldsymbol{\Omega} \bJ = \bJ_1^T \bJ_2^T \left( \bJ_3^T \boldsymbol{\Omega} \bJ_3 \right) \bJ_2 \bJ_1 = \bJ_1^T \left( \bJ_2^T \boldsymbol{\Omega} \bJ_2 \right) \bJ_1 = \bJ_1^T \boldsymbol{\Omega} \bJ_1 = \boldsymbol{\Omega}.
\end{align*}
Each shear has unit determinant: $\det(\bJ_1) = \det(\bJ_2) = \det(\bJ_3) = 1$.
Therefore
\begin{align*}
    \det(\bJ) = \det(\bJ_3)\det(\bJ_2)\det(\bJ_1) = 1 \cdot 1 \cdot 1 = 1.
\end{align*}
}

\subsubsection{Exercises}

\addExcr{%
\begin{enumerate}
    \item \exDerive \textbf{The discrete Gronwall lemma and Forward Euler convergence.}\enspace
    \begin{enumerate}
        \item Let $\by(t)$ solve $\by' = \boldsymbol{f}(t, \by)$ with $\|\boldsymbol{f}(t, \bu) - \boldsymbol{f}(t, \bv)\| \leq L\|\bu - \bv\|$. Use Taylor's theorem to show that the one-step defect of Forward Euler satisfies $\|\by(t_{n+1}) - (\by(t_n) + h\boldsymbol{f}(t_n, \by(t_n)))\| \leq \frac{h^2}{2} M_2$, where $M_2 = \max \|\by''(t)\|$.
        \item Derive the error recurrence $\|\be_{n+1}\| \leq (1 + Lh)\|\be_n\| + \frac{h^2}{2}M_2$.
        \item Apply Lem~\ref{lem:discrete-gronwall} to establish the rigorous global error bound $\|\be_n\| \leq \frac{M_2}{2L}(e^{L(t_n - t_0)} - 1)h$.
    \end{enumerate}
    \item \exDerive \textbf{A-stability and L-stability of implicit methods.}\enspace
    \begin{enumerate}
        \item Plot the boundary of the Forward Euler stability region $|1 + z| = 1$ in the complex plane. Show that for real $\lambda < 0$, stability requires $h \leq 2/|\lambda|$.
        \item For the implicit Trapezoidal Rule, prove that $|R_{\mathrm{TR}}(z)| = 1$ if and only if $\operatorname{Re}(z) = 0$. Explain why this implies that undamped physical oscillations are integrated without artificial numerical dissipation.
        \item Apply the Trapezoidal Rule and Backward Euler to the stiff problem $y' = -1000y$, $y(0) = 1$ with $h = 0.1$. Show that Trapezoidal Rule produces oscillatory values $(-99/101)^n \approx (-0.98)^n$, whereas Backward Euler dampens the mode immediately to $(1/101)^n \approx 10^{-2n}$.
    \end{enumerate}
    \item \exExplore \textbf{Symplectic geometry of the harmonic oscillator.}\enspace Consider the harmonic oscillator $q' = v$, $v' = -q$ with exact energy $E(q, v) = \frac{1}{2}(q^2 + v^2)$.
    \begin{enumerate}
        \item Compute the one-step update matrix $\mathbf{M}_{\mathrm{FE}}$ for Forward Euler and show that $\det(\mathbf{M}_{\mathrm{FE}}) = 1 + h^2 > 1$.
        \item Compute the one-step update matrix $\mathbf{M}_{\mathrm{Verlet}}$ for Velocity Verlet. Confirm that $\det(\mathbf{M}_{\mathrm{Verlet}}) = 1$.
        \item Show that Velocity Verlet conserves an exact modified shadow Hamiltonian $\tilde{H}(q, v) = \frac{1}{2}(q^2 + v^2) - \frac{h^2}{4}q^2 = H(q, v) + O(h^2)$.
    \end{enumerate}
    \item \exCompute \textbf{Adaptive step size control via embedded Runge--Kutta pairs.}\enspace
    \begin{enumerate}
        \item An embedded pair computes a lower-order approximation $\by_{n+1}^{[p]}$ and a higher-order approximation $\by_{n+1}^{[p+1]}$ sharing stages. Show that the difference $\be_{n+1} = \by_{n+1}^{[p+1]} - \by_{n+1}^{[p]}$ satisfies $\|\be_{n+1}\| = C h^{p+1} + O(h^{p+2})$.
        \item Given error tolerance $\varepsilon$ and current step size $h_{\mathrm{old}}$, derive the optimal step size formula $h_{\mathrm{new}} = s h_{\mathrm{old}} (\varepsilon / \|\be_{n+1}\|)^{1/(p+1)}$ with safety factor $s \in (0, 1)$.
        \item Explain why adaptive solvers reject a step and repeat with smaller $h$ when $\|\be_{n+1}\| > \varepsilon$, and why step growth is capped by $h_{\mathrm{new}} \leq 5 h_{\mathrm{old}}$ in production codes.
    \end{enumerate}
    \item \exExplore \textbf{Stiffness in chemical cascades and BDF methods.}\enspace Consider the singularly perturbed initial value problem $\varepsilon y' = -y + \cos(t)$ with $y(0) = 0$ and $\varepsilon = 10^{-4}$.
    \begin{enumerate}
        \item Find the exact solution and identify the fast boundary layer transient $e^{-t/\varepsilon}$ and the slow quasi-steady state $\cos(t) + \varepsilon \sin(t)$.
        \item Explain why explicit RK4 requires $h \lesssim 2.8 \times 10^{-4}$ throughout the entire simulation, even after the transient has decayed to zero.
        \item Formulate the BDF2 update $\frac{3\by_{n+1} - 4\by_n + \by_{n-1}}{2h} = \boldsymbol{f}(t_{n+1}, \by_{n+1})$ and show that it allows time steps $h = 0.1$ in the smooth regime.
    \end{enumerate}
\end{enumerate}
}
\label{ex:ode-exercises}

\addReading{%
Hairer, N{\o}rsett, and Wanner give a careful account of ODE methods and their analysis \citep{hairer1993}. Dahlquist's paper is the classical source for the stability barrier for linear multistep methods \citep{dahlquist1963}. Quarteroni, Sacco, and Saleri provide a broader numerical-analysis reference \citep{quarteroni2007}.
}

\end{document}
